\chapter{Cardiac Ablation}
\index{Cardiac Ablation}

\section*{Definition and Overview}
Cardiac ablation is a procedure used to treat heart rhythm problems (arrhythmias) by deliberately damaging small, precisely targeted areas of heart tissue that are generating or conducting abnormal electrical signals. A thin, flexible tube called a catheter is guided through a blood vessel, usually from the groin, into the heart. Energy such as radiofrequency heat or extreme cold is applied to the target tissue to create a small scar that blocks the faulty electrical pathway. The aim is to restore and maintain a normal heart rhythm without the need for open-heart surgery. The procedure is performed by specialist cardiologists known as electrophysiologists and is usually carried out in a specialised cardiac catheter laboratory.

\section*{Epidemiology}
Heart rhythm disorders are common, particularly in older adults. Atrial fibrillation, the arrhythmia most often treated with ablation, affects roughly one in thirty adults and becomes substantially more common after the age of 65. Supraventricular tachycardias, atrial flutter, and ventricular arrhythmias also affect significant numbers of people. With increasingly effective and safer techniques, the number of ablation procedures performed worldwide has grown steadily over recent decades, and the procedure is now offered to many patients who previously relied on long-term medication. It is one of the most commonly performed cardiac procedures in developed countries.

\section*{Aetiology and Pathophysiology}
The heart normally beats because electrical impulses arise in the sinoatrial node and spread through specialised conducting pathways, producing an orderly contraction that pumps blood efficiently. Arrhythmias occur when this system is disrupted:
\begin{itemize}
    \item \textbf{Abnormal pacemaker tissue:} areas of tissue that fire electrical impulses spontaneously and override the heart's normal pacemaker
    \item \textbf{Re-entry circuits:} extra or damaged electrical pathways that allow impulses to circulate repeatedly, causing rapid heart rates
    \item \textbf{Trigger foci:} single abnormal cells or small patches that fire prematurely, commonly located at the openings of the pulmonary veins in atrial fibrillation
    \item \textbf{Structural heart disease:} scarring from a previous heart attack, valve disease, or cardiomyopathy that disturbs electrical conduction
    \item \textbf{Genetic and metabolic factors:} inherited channel disorders, thyroid disease, or electrolyte imbalance
\end{itemize}

\section*{Clinical Features}
Symptoms depend on the type of arrhythmia and how fast or irregular the heart is beating:
\begin{itemize}
    \item Palpitations -- a sensation of fluttering, racing, or skipping heartbeats
    \item Light-headedness or fainting (syncope)
    \item Shortness of breath, especially on exertion
    \item Chest discomfort or pressure
    \item Fatigue and reduced exercise tolerance
    \item In some cases, no symptoms at all, with the arrhythmia discovered on routine examination or an ECG for another reason
\end{itemize}

\section*{Differential Diagnosis}
The symptoms of an arrhythmia overlap with many other conditions, and ablation is only appropriate when the cause has been identified as an abnormal heart rhythm:
\begin{itemize}
    \item Anxiety, panic attacks, or stress-related palpitations
    \item Anaemia, thyroid disease, or fever causing a fast heart rate
    \item Heart valve disease or cardiomyopathy
    \item Coronary artery disease presenting with chest pain or breathlessness
    \item Premature atrial or ventricular beats that are benign and do not require ablation
    \item Effects of caffeine, alcohol, or recreational drugs on the heart rhythm
\end{itemize}

\section*{When to Seek Medical Care}
Seek medical advice if palpitations are frequent, prolonged, or accompanied by dizziness, fainting, or breathlessness. An unexplained fast or irregular heartbeat should always be assessed.
\textbf{Contact your local emergency services for:}
\begin{itemize}
    \item Sustained chest pain or pressure
    \item Fainting, collapse, or loss of consciousness
    \item Severe breathlessness at rest
    \item A rapid heartbeat that does not settle, with light-headedness or chest pain
\end{itemize}

\section*{Diagnosis and Investigations}
Before ablation, the arrhythmia must be accurately characterised:
\begin{itemize}
    \item \textbf{History and examination:} details of the symptoms, triggers, and any underlying heart disease
    \item \textbf{Electrocardiogram (ECG):} records the heart's electrical activity and captures the arrhythmia if present at the time
    \item \textbf{Holter or event monitor:} portable devices worn for days to weeks to capture intermittent arrhythmias
    \item \textbf{Echocardiogram:} ultrasound to assess heart structure, valve function, and pumping performance
    \item \textbf{Exercise stress testing:} to provoke rhythm problems related to exertion
    \item \textbf{Electrophysiology study:} catheter-based mapping of the heart's electrical circuits, often performed at the start of the ablation itself
    \item \textbf{Blood tests:} thyroid function, electrolytes, and kidney function before the procedure
\end{itemize}

\section*{Management}
Ablation is one treatment option and is chosen after weighing its risks and benefits against long-term medication:
\begin{itemize}
    \item \textbf{Medicines first:} rate-control or rhythm-control drugs are often tried, especially when the arrhythmia is intermittent or mild
    \item \textbf{Ablation procedure:} performed under sedation or general anaesthesia; energy is applied to the mapped abnormal tissue
    \item \textbf{Recovery:} typically one night in hospital; a small dressing at the groin site; most people return to normal activity within one to two weeks
    \item \textbf{Anticoagulation:} blood-thinning medication is often continued or started to reduce the risk of stroke, especially in atrial fibrillation
    \item \textbf{Repeat ablation:} some arrhythmias recur and may require a second procedure
    \item \textbf{Follow-up:} regular review with ECG and monitoring to confirm rhythm control and detect recurrence
\end{itemize}

\section*{Complications}
Most ablations are uncomplicated, but the following can occur:
\begin{itemize}
    \item Bruising, bleeding, or infection at the catheter entry site
    \item Damage to a blood vessel or to the heart wall, occasionally causing bleeding around the heart
    \item Disruption of the heart's normal pacemaker, occasionally requiring a permanent pacemaker
    \item Stroke or transient ischaemic attack from loosened clots or air bubbles
    \item Narrowing of a pulmonary vein as a late complication of atrial fibrillation ablation
    \item Rarely, damage to the oesophagus, diaphragm, or a heart valve
    \item Recurrence of the arrhythmia despite initially successful treatment
\end{itemize}

\section*{Prognosis and Outcomes}
In many cases, ablation markedly reduces or eliminates symptoms of palpitations and improves quality of life. Success rates are highest for regular arrhythmias such as supraventricular tachycardia, where the procedure cures the problem in most people. Results in atrial fibrillation are good but less certain, and some patients require more than one procedure or ongoing medication. Outcomes depend on the type of arrhythmia, the health of the underlying heart, and the experience of the treating centre, so individual expectations should be discussed with the cardiologist before the procedure.

\section*{Prevention and Self-Care}
\begin{itemize}
    \item Limit or avoid alcohol, caffeine, and stimulant drugs that can trigger palpitations
    \item Manage stress and prioritise adequate sleep
    \item Keep blood pressure, cholesterol, and blood sugar well controlled
    \item Maintain a healthy weight and stay physically active within your doctor's advice
    \item Take prescribed rhythm and blood-thinning medicines exactly as directed
    \item Report any new or recurrent palpitations promptly rather than waiting for them to worsen
\end{itemize}

\section*{Sources and Further Reading}
The following sources provide authoritative, up-to-date information on cardiac ablation and its role in treating arrhythmias.
\begin{itemize}
    \item NHS. \textit{Catheter ablation for an irregular heartbeat.} https://www.nhs.uk/conditions/
    \item Mayo Clinic. \textit{Cardiac ablation: what it involves.} https://www.mayoclinic.org/tests-procedures/cardiac-ablation/
    \item National Heart, Lung, and Blood Institute. \textit{Heart arrhythmia and catheter ablation.} https://www.nhlbi.nih.gov/health/heart-arrhythmia
    \item British Heart Foundation. \textit{Catheter ablation for heart rhythm problems.} https://www.bhf.org.uk/
    \item MedlinePlus. \textit{Catheter ablation for heart arrhythmia.} https://medlineplus.gov/arrhythmia.html
\end{itemize}
